\documentclass[
    language=english,
    doctype=lesson-plan,
    institution=none,
]{../../omnilatex}

\RequirePackage{config/document-settings}

\begin{document}

\maketitle

\section*{Lesson Overview}

\begin{description}
    \item[Subject] Computer Science --- Distributed Systems
    \item[Topic] Consensus Protocols: Raft vs Paxos
    \item[Duration] 75 minutes
    \item[Level] Upper undergraduate / graduate
    \item[Prerequisites] Basic understanding of distributed systems,
      network partitions, leader election
\end{description}

\section*{Learning Objectives}
By the end of this lesson, students will be able to:
\begin{enumerate}
    \item Explain the Raft consensus protocol step by step
    \item Compare Raft and Paxos on understandability and performance
    \item Implement a simplified leader election in pseudocode
    \item Identify failure modes and recovery strategies
\end{enumerate}

\section*{Materials}
\begin{itemize}
    \item Projector and slides (beamer presentation)
    \item Printed handouts of Raft paper (Sections 5--8)
    \item Whiteboard markers
    \item Student laptops (for interactive simulation)
\end{itemize}

\section*{Lesson Plan}

\subsection{Warm-Up (10 minutes)}
\textbf{Activity:} Ask students to recall the FLP impossibility result
and discuss why practical systems still achieve consensus.

\begin{itemize}
    \item Poll: ``Can you solve consensus in an asynchronous system?'' (Yes/No/Depends)
    \item Brief discussion of synchrony assumptions
\end{itemize}

\subsection{Direct Instruction (25 minutes)}
\textbf{Content:} Walk through the Raft protocol:

\begin{enumerate}
    \item Leader election (terms, RequestVote RPC)
    \item Log replication (AppendEntries RPC)
    \item Safety (election restriction, commitment)
\end{enumerate}

\textbf{Visual aid:} Draw the Raft state machine on the whiteboard.
Use the interactive simulation at raft.github.io to demonstrate
leader election under network partitions.

\subsection{Guided Practice (20 minutes)}
\textbf{Activity:} Students work in pairs to trace through a Raft
execution scenario:

\begin{enumerate}
    \item 5-node cluster, node 3 is leader
    \item Node 3 crashes
    \item Nodes 1, 2, 4, 5 hold election
    \item Node 4 wins election
    \item Node 3 recovers and rejoins
\end{enumerate}

\textbf{Deliverable:} Each pair produces a sequence diagram showing
the message exchanges.

\subsection{Independent Practice (15 minutes)}
\textbf{Task:} Write pseudocode for the RequestVote RPC handler.
Consider:
\begin{itemize}
    \item When should a node vote?
    \item What if the candidate's log is stale?
    \item How does the term number change?
\end{itemize}

\subsection{Closure (5 minutes)}
\textbf{Exit ticket:} Students write one thing they learned and one
question they still have.

\section*{Assessment}
\begin{itemize}
    \item Observation during pair work
    \item Review of sequence diagrams
    \item Exit ticket responses
\end{itemize}

\section*{Homework}
Read Raft paper Section 9 (Formal Specification) and write a
one-page comparison of Raft's safety proof approach vs Paxos.

\section*{Differentiation}
\begin{description}
    \item[For advanced students] Explore multi-Raft (sharded) systems
    \item[For struggling students] Focus on leader election only;
      skip log replication details
\end{description}

\end{document}
